\MathBookSourcePage{106}
\setcounter{statement}{6}
\begin{Remark}\label{rem:stokes-vector-potential-multiply-connected}
速度向量 $\mathbf{u}$ 只有在 $\Omega$ 单连通时才有满足
\eqref{eq:stokes-vector-potential-tangential-boundary} 的向量势.
不过值得指出, 即使 $\Omega$ 多连通, 问题
\eqref{eq:stokes-vector-biharmonic-tangential} 仍然适定.
事实上, 容易证明下面的范数等价性.
\end{Remark}

\setcounter{statement}{1}
\begin{Lemma}\label{lem:stokes-vector-biharmonic-norm-tangential}
除 Theorem \ref{thm:stokes-dirichlet-well-posedness} 的假设外,
再假设 $\Gamma$ 属于 $\mathcal{C}^{1,1}$ 类或者 $\Omega$ 单连通.
则映射
\[
\boldsymbol{\psi}\longmapsto
\|\Delta\boldsymbol{\psi}\|_{0,\Omega}
\]
是 $\Psi$ 上与 $\|\boldsymbol{\psi}\|$ 等价的一个范数.
\end{Lemma}
\begin{Proof}
首先, 由 Theorem \ref{thm:poincare-friedrichs},
$\|\boldsymbol{\psi}\|$ 在 $\Psi$ 上等价于
\[
\left\{
\|\boldsymbol{\psi}\|_{0,\Omega}^2
+\|\operatorname{div}\boldsymbol{\psi}\|_{1,\Omega}^2
+|\operatorname{curl}\boldsymbol{\psi}|_{1,\Omega}^2
\right\}^{1/2}.
\]
其次, 注意
\eqref{eq:stokes-vector-laplacian-decomposition}
正是向量 $\Delta\boldsymbol{\psi}$ 的正交分解. 这一正交性给出
\[
\|\Delta\boldsymbol{\psi}\|_{0,\Omega}^2
=
\|\operatorname{curl}\operatorname{curl}
\boldsymbol{\psi}\|_{0,\Omega}^2
+|\operatorname{div}\boldsymbol{\psi}|_{1,\Omega}^2.
\]
但由于
$\operatorname{curl}\boldsymbol{\psi}\in H_0^1(\Omega)^3$,
由 Remark \ref{rem:sec2-hdiv-hcurl-estimate} 推得
\[
|\operatorname{curl}\boldsymbol{\psi}|_{1,\Omega}
\cong
\|\operatorname{curl}\operatorname{curl}
\boldsymbol{\psi}\|_{0,\Omega}.
\]
此外,
$\int_\Omega\operatorname{div}\boldsymbol{\psi}\,\mathrm{d}x=0$
意味着
\[
\|\operatorname{div}\boldsymbol{\psi}\|_{1,\Omega}
\cong
|\operatorname{div}\boldsymbol{\psi}|_{1,\Omega}.
\]
事实上, 直接应用 Theorem
\ref{thm:peetre-tartar-compact-operator} 的 $3^\circ)$ 可得
\[
\|v\|_{1,\Omega}
\cong
\left(
|v|_{1,\Omega}^2
+\left|\int_\Omega v\,\mathrm{d}x\right|^2
\right)^{1/2}
\qquad
\forall v\in H^1(\Omega).
\]
所以只需证明
\begin{equation}\label{eq:stokes-vector-potential-poincare}
\|\boldsymbol{\psi}\|_{0,\Omega}
\le
C\left\{
\|\operatorname{div}\boldsymbol{\psi}\|_{0,\Omega}^2
+\|\operatorname{curl}\boldsymbol{\psi}\|_{0,\Omega}^2
\right\}^{1/2}
\qquad
\forall\boldsymbol{\psi}\in\Psi.
\end{equation}
当 $\Omega$ 单连通时, 这由 Lemma
\ref{lem:sec3-w0-curl-isomorphism} 证明.
事实上, 即使 $\Gamma$ 有多个连通分支, 其论证仍说明
\[
\|\boldsymbol{\psi}\|_{0,\Omega}
\le
C\|\operatorname{curl}\boldsymbol{\psi}\|_{0,\Omega}
\quad
\forall\boldsymbol{\psi}\in\Psi
\quad\text{with}\quad
\operatorname{div}\boldsymbol{\psi}=0.
\]
若 $\operatorname{div}\boldsymbol{\psi}\ne0$, 总可以考虑差
$\boldsymbol{\psi}-\mathbf{w}$, 得到 $\Psi$ 中的无散函数,
其中 $\mathbf{w}\in H_0^1(\Omega)^3$ 满足(参见 Corollary
\ref{cor:sec2-grad-div-isomorphisms})
\[
\operatorname{div}\mathbf{w}
=\operatorname{div}\boldsymbol{\psi},
\qquad
|\mathbf{w}|_{1,\Omega}
\le
C\|\operatorname{div}\boldsymbol{\psi}\|_{0,\Omega}.
\]
当 $\Gamma$ 属于 $\mathcal{C}^{1,1}$ 类时,
\eqref{eq:stokes-vector-potential-poincare} 容易由 Theorem
\ref{thm:sec3-w-h1-embedding}, Theorem
\ref{thm:peetre-tartar-compact-operator} 和 Remark
\ref{rem:sec3-zero-curl-zero-tangential-implies-zero} 得到.
\end{Proof}

\setcounter{statement}{7}
\begin{Remark}\label{rem:stokes-vector-biharmonic-not-original-potential}
还可以证明(参见 Dominguez [24]), 当 $\Omega$ 多连通时,
问题 \eqref{eq:stokes-vector-biharmonic-tangential} 的解不是原 Stokes
问题的势.
\end{Remark}

\MathBookSourcePage{107}
\setcounter{statement}{8}
\begin{Remark}\label{rem:stokes-vector-biharmonic-tangential-strong-form}
Bendali, Dominguez \& Gallic [6] 已经说明, 问题
\eqref{eq:stokes-vector-biharmonic-tangential} 有如下解释:
\begin{equation}\label{eq:stokes-vector-biharmonic-tangential-strong}
\left\{
\begin{aligned}
\nu\Delta^2\boldsymbol{\psi}
&=\operatorname{curl}\mathbf{f}
&&\text{in }H^{-2}(\Omega)^3,\\
\operatorname{div}\boldsymbol{\psi}&=0
&&\text{in }\Omega,\\
(\operatorname{curl}\boldsymbol{\psi})
\times\mathbf{n}|_\Gamma&=\mathbf{0},
\qquad
\boldsymbol{\psi}\times\mathbf{n}|_\Gamma=\mathbf{0},\\
\int_{\Gamma_i}
\boldsymbol{\psi}\cdot\mathbf{n}\,\mathrm{d}s&=0,
&&0\le i\le p.
\end{aligned}
\right.
\end{equation}
\end{Remark}

对满足边界条件
\eqref{eq:stokes-vector-potential-normal-boundary} 的向量势,
可作非常相似的分析. 这里选择空间
\[
\Psi_1
=
\{\,\boldsymbol{\phi}\in L^2(\Omega)^3;
\ \operatorname{div}\boldsymbol{\phi}\in H^1(\Omega),\
\operatorname{curl}\boldsymbol{\phi}\in H_0^1(\Omega)^3,\
\boldsymbol{\phi}\cdot\mathbf{n}|_\Gamma=0\,\},
\]
并赋予它与 $\Psi$ 相同的范数. 于是得到 Lemma
\ref{lem:stokes-vector-biharmonic-tangential} 和 Lemma
\ref{lem:stokes-vector-biharmonic-norm-tangential} 的对应结果.
更确切地, 首先可以证明范数的等价性.

\setcounter{statement}{2}
\begin{Lemma}\label{lem:stokes-vector-biharmonic-norm-normal}
若 $\Omega$ 如 Lemma
\ref{lem:stokes-vector-biharmonic-tangential} 中所述, 则映射
\[
\boldsymbol{\psi}\longmapsto
\|\Delta\boldsymbol{\psi}\|_{0,\Omega}
\]
是 $\Psi_1$ 上与 $\|\boldsymbol{\psi}\|$ 等价的一个范数.
\end{Lemma}

由 Lemma \ref{lem:stokes-vector-biharmonic-norm-normal},
双调和问题

求 $\boldsymbol{\psi}\in\Psi_1$, 使
\begin{equation}\label{eq:stokes-vector-biharmonic-normal}
\nu(\Delta\boldsymbol{\psi},
\Delta\boldsymbol{\phi})
=
\langle\mathbf{f},
\operatorname{curl}\boldsymbol{\phi}\rangle
\qquad
\forall\boldsymbol{\phi}\in\Psi_1
\end{equation}
有唯一解 $\boldsymbol{\psi}\in\Psi_1$.
容易检验 $\boldsymbol{\psi}$ 是 $\mathbf{u}$ 的向量势,
从而 $\operatorname{div}\boldsymbol{\psi}=0$.

\setcounter{statement}{3}
\begin{Lemma}\label{lem:stokes-vector-biharmonic-normal}
假设 $\Omega$ 如 Lemma
\ref{lem:stokes-vector-biharmonic-tangential} 中所述.
双调和问题 \eqref{eq:stokes-vector-biharmonic-normal}
在 $\Psi_1$ 中有唯一解 $\boldsymbol{\psi}$, 并且
\[
\operatorname{div}\boldsymbol{\psi}=0,
\qquad
\operatorname{curl}\boldsymbol{\psi}=\mathbf{u},
\]
其中 $\mathbf{u}$ 是齐次 Stokes 问题
\eqref{eq:stokes-dirichlet-problem} 的解.
\end{Lemma}

\setcounter{statement}{9}
\begin{Remark}\label{rem:stokes-vector-biharmonic-normal-strong-form}
问题 \eqref{eq:stokes-vector-biharmonic-normal} 有如下解释:
\begin{equation}\label{eq:stokes-vector-biharmonic-normal-strong}
\left\{
\begin{aligned}
\nu\Delta^2\boldsymbol{\psi}
&=\operatorname{curl}\mathbf{f}
&&\text{in }H^{-2}(\Omega)^3,\\
\operatorname{div}\boldsymbol{\psi}&=0
&&\text{in }\Omega,\\
\operatorname{curl}\boldsymbol{\psi}|_\Gamma&=\mathbf{0},
\qquad
\boldsymbol{\psi}\cdot\mathbf{n}|_\Gamma=0.
\end{aligned}
\right.
\end{equation}
\end{Remark}

最后考虑非齐次 Stokes 问题. 一方面, 必须假设边界值
$\mathbf{g}$ 满足
\begin{equation}\label{eq:stokes-vector-nonhomogeneous-flux}
\int_{\Gamma_i}\mathbf{g}\cdot\mathbf{n}\,\mathrm{d}s=0,
\qquad 0\le i\le p,
\end{equation}

\MathBookSourcePage{108}
从而 $\mathbf{u}$ 的确有一个向量势 $\boldsymbol{\psi}$.
另一方面, $\boldsymbol{\psi}$ 不能满足边界条件
\eqref{eq:stokes-vector-potential-tangential-boundary},
因为该条件蕴含
$(\operatorname{curl}\boldsymbol{\psi})
\cdot\mathbf{n}|_\Gamma=0$.
但是可以指定条件
\eqref{eq:stokes-vector-potential-normal-boundary}.
因此, 向量势空间的一个合理选择是
\[
\Psi_2
=
\{\,\boldsymbol{\phi}\in L^2(\Omega)^3;
\ \operatorname{div}\boldsymbol{\phi}\in H^1(\Omega),\
\operatorname{curl}\boldsymbol{\phi}\in H^1(\Omega)^3,\
\boldsymbol{\phi}\cdot\mathbf{n}|_\Gamma=0\,\}.
\]
同样, 取 $\Psi_1$ 为试验函数空间. 考虑双调和问题:

求 $\boldsymbol{\psi}\in\Psi_2$, 且
$\operatorname{curl}\boldsymbol{\psi}|_\Gamma=\mathbf{g}$, 使
\begin{equation}\label{eq:stokes-vector-biharmonic-nonhomogeneous}
\nu(\Delta\boldsymbol{\psi},
\Delta\boldsymbol{\phi})
=
\langle\mathbf{f},
\operatorname{curl}\boldsymbol{\phi}\rangle
\qquad
\forall\boldsymbol{\phi}\in\Psi_1.
\end{equation}
容易检验这个问题有唯一解, 但
$\boldsymbol{\psi}$ 与 $\mathbf{u}$ 的关系并非完全显然,
因为 \eqref{eq:stokes-vector-curlcurl-graddiv-orthogonality}
不再对所有 $\boldsymbol{\phi}\in\Psi_1$ 成立.
不过, 注意 $\boldsymbol{\psi}$ 和 $\mathbf{u}$ 在 $\Psi_1$ 中的无散势
$\boldsymbol{\psi}_0$,
\[
\operatorname{curl}\boldsymbol{\psi}_0=\mathbf{u},
\qquad
\operatorname{div}\boldsymbol{\psi}_0=0,
\qquad
\boldsymbol{\psi}_0\cdot\mathbf{n}|_\Gamma=0,
\]
都满足
\[
\nu(
\operatorname{curl}\operatorname{curl}\boldsymbol{\psi},
\operatorname{curl}\operatorname{curl}\boldsymbol{\phi})
=
\langle\mathbf{f},
\operatorname{curl}\boldsymbol{\phi}\rangle
\]
对所有满足
$\operatorname{div}\boldsymbol{\phi}=0$ 的
$\boldsymbol{\phi}\in\Psi_1$ 成立. 因此置
\[
\mathbf{w}
=
\operatorname{curl}
(\boldsymbol{\psi}-\boldsymbol{\psi}_0)
\in V,
\]
便得到
$\operatorname{curl}\mathbf{w}=\mathbf{0}$, 从而
$\mathbf{w}=\mathbf{0}$. 因此
$\boldsymbol{\psi}$ 是 $\mathbf{u}$ 的向量势,
但它不一定无散.

\setcounter{statement}{5}
\begin{Theorem}\label{thm:stokes-vector-biharmonic-nonhomogeneous}
设 $\Omega$ 如 Lemma
\ref{lem:stokes-vector-biharmonic-tangential} 中所述, 并设
$\mathbf{g}$ 满足 \eqref{eq:stokes-vector-nonhomogeneous-flux}.
双调和问题
\eqref{eq:stokes-vector-biharmonic-nonhomogeneous}
在 $\Psi_2$ 中有唯一解 $\boldsymbol{\psi}$, 且
$\operatorname{curl}\boldsymbol{\psi}$ 是非齐次 Stokes 问题
\eqref{eq:stokes-dirichlet-problem} 的唯一解.
\end{Theorem}

\setcounter{statement}{10}
\begin{Remark}\label{rem:stokes-vector-potential-divergence-neumann}
$\boldsymbol{\psi}$ 的散度是 Neumann 问题
\[
\Delta\lambda=0\quad\text{in }\Omega,
\qquad
\int_\Omega\lambda\,\mathrm{d}x=0,
\qquad
\frac{\partial\lambda}{\partial n}
=(\operatorname{curl}\mathbf{u})\cdot\mathbf{n}
\quad\text{on }\Gamma
\]
的解 $\lambda$. 因此
$\operatorname{div}\boldsymbol{\psi}=0$ 当且仅当
$(\operatorname{curl}\mathbf{u})\cdot\mathbf{n}|_\Gamma=0$.
\end{Remark}

\setcounter{statement}{11}
\begin{Remark}\label{rem:stokes-vector-potential-divergence-boundary-constraint}
若希望某个双调和问题的解是 $\boldsymbol{\psi}_0$,
则必须在空间 $\Psi_1$ 和 $\Psi_2$ 中加入约束
$(\operatorname{div}\boldsymbol{\phi})|_\Gamma=0$.
\end{Remark}

\setcounter{statement}{12}
\begin{Remark}\label{rem:stokes-vector-potential-surface-conditions}
还可以通过一方面置
\[
(\operatorname{curl}\boldsymbol{\psi})\times\mathbf{n}
=\mathbf{g}\times\mathbf{n}
\quad\text{on }\Gamma,
\]
另一方面置

\MathBookSourcePage{109}
\[
\operatorname{div}_\Gamma
(\boldsymbol{\psi}\times\mathbf{n})
=\mathbf{g}\cdot\mathbf{n},
\qquad
\operatorname{curl}_\Gamma
(\boldsymbol{\psi}\times\mathbf{n})=0,
\]
来指定条件
$(\operatorname{curl}\boldsymbol{\psi})|_\Gamma=\mathbf{g}$.
下标 $\Gamma$ 表明算子
$\operatorname{div}$ 和 $\operatorname{curl}$ 是曲面算子.
更多细节参见 Roux [69].
\end{Remark}
